\documentclass[10pt,twoside,titlepage]{article}   % Standard LaTeX
\usepackage{amsmath,amssymb}    % For better support of math
\usepackage{amsthm}
\usepackage{url}
\usepackage{graphicx}          % Enable for eps,jpeg figures

\usepackage{afterpage}
\usepackage{float}

\newcommand\blankpage{%
    \null
    \thispagestyle{empty}%
    \addtocounter{page}{-1}%
    \newpage}

\urldef\HHIURL\url{https://en.wikipedia.org/wiki/Herfindahl%E2%80%93Hirschman_index#Effective_assets_in_a_portfolio}
\urldef\CORRURL\url{https://stats.stackexchange.com/questions/226380/derivation-of-the-standard-error-for-pearsons-correlation-coefficient}
\usepackage{dsfont}

\DeclareMathOperator*{\argmax}{arg\,max}
\DeclareMathOperator*{\argmin}{arg\,min}

%% Controls spacing between lines (\doublespacing, \onehalfspacing, etc.):
\usepackage{setspace}

%% Use the line below for official NYU version, which requires
%% double line spacing. For all other uses, this is unnecessary,
%% so the line can be commented out.
\doublespacing % requires package setspace, invoked above

\usepackage{listings}
\usepackage{color}

\definecolor{dkgreen}{rgb}{0,0.6,0}
\definecolor{gray}{rgb}{0.5,0.5,0.5}
\definecolor{mauve}{rgb}{0.58,0,0.82}

\begin{document}

\title{Extending the Oracle Approximating Shrinkage to General Distribution}
\author{Michael J. Lewis}
\date{\copyright\today}         % change \today to a particular date

\maketitle

\begin{abstract}
The Oracle Approximating Shrinkage (OAS) was created assuming an underlying Gaussian distribution.
In this work, we extend the tool to all distributions with 4th moments.
\end{abstract}

\section{The formulation}\label{sec-formulation}
This work should be seen as an extension of \cite{OAS}. To set the stage, we first port over the relevant equations.
\begin{equation}
\rho_O = \frac{ E \left[ \operatorname{Tr} \left( \left( \Sigma - \hat{S} \right) \left( \hat{F} - \hat{S} \right) \right) \right] }{ E \left[ \lVert \hat{S} - \hat{F} \rVert_F^2 \right] }
\end{equation}
\begin{equation}
E \left[ \operatorname{Tr} \left( \left( \Sigma - \hat{S} \right) \left( \hat{F} - \hat{S} \right) \right) \right] =
\frac{\operatorname{Tr}( \Sigma )}{p} E \left[ \operatorname{Tr}(\hat{S})\right] 
- \frac{ E \left[ \operatorname{Tr}^2(\hat{S})\right] }{ p } \\
- E \left[ \operatorname{Tr}(\Sigma \hat{S})\right]
+ E \left[ \operatorname{Tr}(\hat{S}^2)\right]
\end{equation}
\begin{equation}
E \left[ \lVert \hat{S} - \hat{F} \rVert_F^2 \right] =
E \left[ \operatorname{Tr}(\hat{S}^2)\right]
- \frac{ E \left[ \operatorname{Tr}^2(\hat{S})\right] }{ p }
\end{equation}
\begin{equation}
E \left[ \operatorname{Tr}(\hat{S})\right] = \operatorname{Tr}(\Sigma)
\end{equation}
It is at this point, we will deviate slightly from \cite{OAS} by considering the general distribution case.
Observe that
\begin{equation*}
\begin{aligned}
E \left[ \operatorname{Tr}(\hat{S}^2)\right] &= E \left[ \lVert \hat{S} \rVert_F^2 \right] \\
&= \sum_{i,j} \frac{1}{n^2} \sum_{k,l} E \left[ X_i^k X_j^k X_i^l X_j^l \right] \\
&= \sum_{i,j} \frac{1}{n^2} \left( \sum_{k \neq l} E \left[ X_i^k X_j^k X_i^l X_j^l \right] + \sum_{k = l} E \left[ X_i^k X_j^k X_i^l X_j^l \right] \right) \\
&= \frac{n-1}{n} \sum_{i,j} \Sigma^2_{i,j} + \frac{1}{n^2} \sum_{i,j} \sum_k E \left[ (X_i^k)^2 (X_j^k)^2 \right] \\
&= \frac{n-1}{n} \operatorname{Tr}(\Sigma^2) + \frac{1}{n} \sum_{i,j} E \left[ X_i^2 X_j^2 \right] \\
\end{aligned}
\end{equation*}
Similarly,
\begin{equation*}
\begin{aligned}
E \left[ \operatorname{Tr^2}(\hat{S})\right] &= E \left[ \sum_{i,j} S_{ii} S_{jj} \right] \\
&= \sum_{i,j} \frac{1}{n^2} \sum_{k,l} E \left[  (X_i^k)^2  (X_j^l)^2  \right] \\
&= \sum_{i,j} \frac{1}{n^2} \left( \sum_{k \neq l} E \left[  (X_i^k)^2  (X_j^l)^2  \right] + \sum_{k = l} E \left[  (X_i^k)^2  (X_j^l)^2  \right] \right) \\
&= \frac{n-1}{n} \sum_{i,j} \Sigma_{ii} \Sigma_{jj} + \frac{1}{n^2} \sum_{i,j} \sum_k E \left[ (X_i^k)^2  (X_j^k)^2 \right] \\
&= \frac{n-1}{n} \operatorname{Tr^2}(\Sigma) + \frac{1}{n} \sum_{i,j} E \left[ X_i^2  X_j^2 \right] \\
\end{aligned}
\end{equation*}
In both cases, the term $E \left[ X_i^2  X_j^2 \right]$ shows up.
For a centered general distribution with a 4th moment, we have
\begin{equation}
E \left[ X_i^2  X_j^2 \right] = \Sigma_{ii} \Sigma_{jj} + 2 \Sigma_{ij}^2 + \kappa( X_i, X_i, X_j, X_j )
\end{equation}
where $\kappa( X_i, X_i, X_j, X_j )$ represents the fourth-order joint cumulant. 
Note that this cumulant vanishes for a gaussian distribution. 
Thus,
\begin{equation}
\label{fourth_moment_eqn}
\sum_{i,j} E \left[ X_i^2  X_j^2 \right] = \operatorname{Tr^2}(\Sigma) + 2 \operatorname{Tr}(\Sigma^2) + \sum_{i,j} \kappa( X_i, X_i, X_j, X_j )
\end{equation}
Let $\alpha =\sum_{i,j} \kappa( X_i, X_i, X_j, X_j )$. 

Returning to the above,
\begin{equation}
E \left[ \operatorname{Tr}(\hat{S}^2)\right] = \frac{n+1}{n} \operatorname{Tr}(\Sigma^2) + \frac{1}{n} \operatorname{Tr^2}(\Sigma) + \frac{\alpha}{n}
\end{equation}
\begin{equation}
E \left[ \operatorname{Tr^2}(\hat{S})\right] = \operatorname{Tr^2}(\Sigma) + \frac{2}{n} \operatorname{Tr}(\Sigma^2) + \frac{\alpha}{n}
\end{equation}
which is very similar to the results in \cite{OAS} with the extra $\frac{\alpha}{n}$ term.
Collating the results above gives us
\begin{equation}
\label{opt_rho}
\rho_O = \frac{ (1-2/p) \operatorname{Tr}(\Sigma^2) + \operatorname{Tr^2}(\Sigma) + \alpha (1 - 1/p) }{ (n+1-2/p) \operatorname{Tr}(\Sigma^2) + (1 - n/p) \operatorname{Tr^2}(\Sigma) + \alpha (1 - 1/p) }
\end{equation}

\section{Extending OAS}\label{sec-extension}
With the relevant formulation out of the way, we now consider the original design of OAS, but extended to include our new term $\alpha (1 - 1/p)$.
As in \cite{OAS}, we have
\begin{equation}
\label{rho_eval}
\hat{\rho}_{j+1} = \frac{ (1-2/p) \operatorname{Tr}(\hat{\Sigma}_j \hat{S} ) + \operatorname{Tr^2}( \hat{\Sigma}_j ) + \alpha (1 - 1/p) }{ (n+1-2/p) \operatorname{Tr}( \hat{\Sigma}_j \hat{S} ) + (1 - n/p) \operatorname{Tr^2}( \hat{\Sigma}_j ) + \alpha (1 - 1/p) } \\
\end{equation}
\begin{equation}
\label{sigma_eval}
\hat{\Sigma}_{j+1} = (1 - \hat{\rho}_{j+1}) \hat{S} + \hat{\rho}_{j+1} \hat{F} \\
\end{equation}
which leverages \eqref{opt_rho}, but replaces $\operatorname{Tr}(\Sigma)$ and $\operatorname{Tr}(\Sigma^2)$ with 
$\operatorname{Tr}(\hat{\Sigma}_j)$ and $\operatorname{Tr}(\hat{\Sigma}_j \hat{S})$, respectively.
Let
\begin{equation}
\hat{\phi} = \frac{ \operatorname{Tr}( \hat{S}^2 ) - \operatorname{Tr^2}( \hat{S} ) / p }{ (1-2/p) \operatorname{Tr}( \hat{S}^2 ) + \operatorname{Tr^2}( \hat{S} ) } = \frac{ \operatorname{Tr}( \hat{S}^2 ) - \operatorname{Tr^2}( \hat{S} ) / p }{ M }
\end{equation}
and $\beta = \alpha (1 -  1/p) / M$.
Plugging in $\Sigma_j$ from \eqref{sigma_eval} and plugging into \eqref{rho_eval} yields
\begin{equation}
\hat{\rho}_{j+1} = \frac{ 1 + \beta - (1 - 2/p) \hat{\phi} \hat{\rho}_j }{ 1 + \beta + n \hat{\phi} - (n+1-2/p) \hat{\phi} \hat{\rho}_j }
\end{equation}
Using the change of variables
\begin{equation}
b_j  = \frac{1}{ 1 + \beta - (n+1-2/p) \hat{\phi} \hat{\rho}_j } \Leftrightarrow \hat{\rho}_j = \frac{ 1 + \beta - \frac{1}{b_j} }{ (n+1-2/p) \hat{\phi} },
\end{equation}
we obtain \eqref{rho_eval} is equivalent to
\begin{equation}
b_{j+1} =  \frac{  n \hat{\phi} }{ 1 + \beta - (1-2/p) \hat{\phi} } b_j + \frac{ 1 }{  1 + \beta - (1-2/p) \hat{\phi} }
\end{equation}
It is easy to see that
\begin{equation}
  \lim_{j \to \infty} b_j =
    \begin{cases}
      \infty & \frac{  n \hat{\phi} }{ 1 + \beta - (1-2/p) \hat{\phi} } \geq  1 \\
      \frac{1}{1 + \beta - (n+1-2/p) \hat{\phi} }      & \frac{  n \hat{\phi} }{ 1 + \beta - (1-2/p) \hat{\phi} } < 1 \\
    \end{cases}
\end{equation}
This implies 
\begin{equation}
  \lim_{j \to \infty} \hat{\rho}_j =
    \begin{cases}
      \frac{ 1 + \beta }{ (n+1-2/p) \hat{\phi} } & (n+1-2/p) \hat{\phi} \geq 1 + \beta \\
      1                                          & (n+1-2/p) \hat{\phi} < 1 + \beta \\
    \end{cases}
\end{equation}
or equivalently, 
\begin{equation}
\label{new_OAS}
\rho_{OAS extended} = \operatorname{min}\left( 1, \frac{ 1 + \beta }{ (n+1-2/p) \hat{\phi} } \right)
\end{equation}

Note that, in the limit, we have $1 + \beta \geq 0$  (The proof is deferred to Appendix~\ref{app:proof}).

What's left is a consistent estimator for the $\beta$ term.
Given that $M = (1-2/p) \operatorname{Tr}( \hat{S}^2 ) + \operatorname{Tr^2}( \hat{S} )$ is already handed, this leaves the $\alpha$ term.
Given its definition and \eqref{fourth_moment_eqn}, the natural choice for the estimator would be
\begin{equation}
\hat{\alpha} = \left( \sum_{i,j} \frac{1}{n} \sum_k (X_i^k)^2 (X_j^k)^2 \right) - \operatorname{Tr^2}(\hat{S}) - 2 \operatorname{Tr}(\hat{S}^2),
\end{equation}
which converges to the true $\alpha$ by the law of large numbers.
However, for finite samples, $\hat{\alpha}$ may drift from $\alpha$, and thus recommend clamping the result, i.e.
\begin{equation}
\hat{\rho}_{OAS extended} = \operatorname{max} \left( 0, \operatorname{min}\left( 1, \frac{ 1 + \beta }{ (n+1-2/p) \hat{\phi} } \right) \right)
\end{equation}

\clearpage

\appendix
\section{Proof That $1 + \beta \geq 0$}
\label{app:proof}
It is obvious from \eqref{new_OAS} that $\rho_{OAS extended} \leq 1$, but leaves open the question of the lower bound.
For this we analyze the numerator $1 + \beta$, and find that it is non-negative.
\begin{proof}
We observe
$\sum_{i,j} E \left[ X_i^2  X_j^2 \right] = E \left[ \lVert X \rVert^4 \right]$ and
$\operatorname{Tr^2}(\Sigma) = \left( E \left[ \lVert X \rVert^2 \right] \right)^2$.
Thus,
\[
\alpha = \sum_{i,j} E \left[ X_i^2  X_j^2 \right] - \operatorname{Tr^2}(\Sigma) - 2 \operatorname{Tr}(\Sigma^2) = \operatorname{Var} \left(  \lVert X \rVert^2 \right) - 2 \operatorname{Tr}(\Sigma^2)
\]
implying $\alpha \geq - 2 \operatorname{Tr}(\Sigma^2)$.
Noting $\beta = \alpha(1-1/p)/M$, it suffices to show
\[
-2\operatorname{Tr}(\Sigma^2) \cdot \frac{1-1/p}{M} \geq -1 \Rightarrow M \geq 2(1-1/p)\operatorname{Tr}(\Sigma^2).
\]
Substituting $M = (1-2/p)\operatorname{Tr}(\Sigma^2) + \operatorname{Tr}^2(\Sigma)$, this reduces to
\[
\operatorname{Tr}^2(\Sigma) \geq \operatorname{Tr}(\Sigma^2).
\]
For any positive semidefinite $\Sigma$ with eigenvalues $\lambda_1, \ldots, \lambda_p \geq 0$,
\[
\operatorname{Tr}^2(\Sigma) - \operatorname{Tr}(\Sigma^2)
= \Bigl(\sum_i \lambda_i\Bigr)^2 - \sum_i \lambda_i^2
= 2\sum_{i < j} \lambda_i \lambda_j \geq 0.
\]
Hence $1 + \beta \geq 0$, and $\rho_{OAS\text{ extended}} \in [0,1]$ for the true $\alpha$.
\end{proof}

\begin{thebibliography}{9}
\bibitem{OAS}
Y. Chen, A. Wiesel, Y. C. Eldar, and A. O. Hero, "Shrinkage algorithms for MMSE covariance estimation," 
IEEE Transactions on Signal Processing, vol. 58, no. 10, pp. 5016–5029, Oct. 2010.
\end{thebibliography}


\end{document}
